% =====================================================================
%  Chapter 30 — The Platform Selection Rule\idx{Selection rule}
% =====================================================================

\chapter{The Platform Selection Rule}
\label{ch:selection-rule}

The book's deliverable is not an architecture. It is a rule that returns an
architecture, and it is indifferent to the answer it gives for any particular
target. A rule that always returns the same platform is a preference. This
chapter writes the rule down, works it on the cases of Chapter 14, and shows
that it declines the lunar platform as often as it selects it.

\section{The inputs}

For a target with a bracketed diameter \(D\), a bracketed density \(\rho\), a
geocentric distance \(\Delta\), a heliocentric distance \(r\), a phase angle
\(\alpha\) and an encounter epoch, the rule runs in six steps.

\begin{enumerate}
\item Compute the required angular resolution. For a resolved shape
measurement with \(N\) samples across the diameter,
\(B \geq N\lambda\Delta/D\), with \(N\) between 4 and 7 and the wavelength
chosen from the target's photometric brightness.
\item Compute the over-resolution limit, which caps the usable baseline at
the point where the body's own structure, rather than the instrument, sets
the visibility.
\item Compute the required astrometric precision as \(\lambda/(B\cdot
\mathrm{SNR})\), and compare it with the along-track budget at the
encounter horizon.
\item Check whether the required baseline is buildable at the candidate
sites: a terrestrial optical array is limited by atmosphere and metrology, a
lunar surface array by construction and logistics, and an orbital formation
by station-keeping and pointing.
\item Check whether the observable can be scheduled inside the target's
brightness window, which for a close approach is a window of days rather
than hours.
\item Check, and only if this passes, whether the physical uncertainty
dominates the encounter probability.
\end{enumerate}

The last condition is the substantive content of the rule.

\section{The condition, in its quantitative form}

The physical term dominates when the along-track displacement produced by
the uncertainty in the Yarkovsky drift exceeds the astrometric along-track
uncertainty at the horizon.

For a 60-metre object at 1.1 astronomical units, a thirty per cent error in
the drift rate produces 17.2 kilometres of along-track error over the seven
years from 2025 to 2032, against a published B-plane uncertainty of 25,000
kilometres for that encounter after a single extra observation, and 270,000
kilometres before one. The astrometric term dominates by three orders of
magnitude at that horizon.

At the 2036 Apophis horizon the situation reverses, because the keyhole
sits in the tail. The diameter assumption alone is worth 507 kilometres of
keyhole displacement and the density assumption 379 kilometres, against an
astrometric standard deviation of 399 kilometres.

\begin{table}[!htbp]
\centering
\small
\caption{The rule applied to the two documented cases.}
\label{tab:rule-applied}
\begin{tabularx}{\textwidth}{@{}P{1.39in}P{1.17in}Y@{}}
\toprule
\rowcolor{panel}
\textbf{Case} & \textbf{Verdict} & \textbf{Why} \\
\midrule
2024 YR4, 2032 & \textbf{decline the lunar platform} & the astrometric
term dominates by three orders of magnitude; characterising it does not
shorten this clock \\
\rowcolor{panel}
Apophis, 2029 & \textbf{select the lunar platform} & the physical terms
are each comparable with the astrometric sigma, and the keyhole is in the
tail \\
\bottomrule
\end{tabularx}
\end{table}

\tabref{tab:rule-applied} is the rule doing its job. For the case that is
the book's worked example in every popular account of 2024 YR4, it declines
the platform. For the case with the largest published impact-probability
sensitivities in the literature, it selects it.

\section{Why the rule is written this way}

Three design decisions in the rule deserve their reasons stated.

The rule is written as a sequence with a conditional at the end, rather than
as a set of weights, because the conditional is a statement about the
physics and not a preference. The last step is a test that can be failed,
and a rule that cannot be failed is a slogan.

The rule returns a platform rather than a measurement. Several platforms can
close the same job --- radar and a lunar array can both deliver a shape ---
and the rule has to be able to return either, on the grounds of cost,
opportunity and what else the platform does.

The rule's cost term is deliberately not a number in this book. The
published cost of a space mission, a ground array and a lunar station are
not commensurable in any way a reader could use without a procurement
context the book does not have. What the rule does instead is require that
the platform be buildable at the required baseline, which is the constraint
that actually discriminates between the candidates: the terrestrial
alternative for a kilometre-scale optical baseline is an atmospheric and
metrological problem, and that is the whole of the discrimination.

\begin{deeper}{What the rule does with the lunar platform's weaknesses}
The rule's fourth step --- buildability\idx{Buildability} at the required baseline --- is where
the lunar case is made, and it is a step the Moon passes because the
construction problem separates from the atmosphere problem. The rule's fifth
step --- schedulability inside the brightness window --- is where the lunar
platform's disadvantages bite, and they are handled by the fact that the
lunar sky rotates slowly enough that a window measured in days is
achievable. The rule does not pretend that dust and the night are absent.
It requires that the campaign structure account for them, which is what the
thirty per cent overhead allowance in Chapter 25 does.
\end{deeper}

\section{The rule as stated for use}

\begin{table}[!htbp]
\centering
\small
\caption{The selection rule, stated for use.}
\label{tab:rule}
\begin{tabularx}{\textwidth}{@{}P{2.57in}Y@{}}
\toprule
\rowcolor{panel}
\textbf{Step} & \textbf{Test} \\
\midrule
1 & Is \(N\lambda\Delta/D\) buildable at any candidate site? If not, no
resolved-shape measurement is possible and the answer is a thermal or radar
measurement. \\
\rowcolor{panel}
2 & Is the matched baseline inside the over-resolution limit? If not, shorten
the baseline or shorten the wavelength. \\
3 & Does \(\lambda/(B\cdot\mathrm{SNR})\) meet the along-track budget at the
encounter horizon? If not, a longer baseline or a longer integration is
required. \\
\rowcolor{panel}
4 & Is the site able to hold the baseline? Terrestrial: no, above a few
hundred metres. Lunar: yes, with a surveyed anchor array. Orbital: station
keeping sets the limit. \\
5 & Is there a brightness window, and is it long enough? \\
\rowcolor{panel}
6 & Does the physical uncertainty exceed the astrometric uncertainty at the
horizon? If not, the answer is more astrometry, and the characterisation
measurement does not matter. \\
\bottomrule
\end{tabularx}
\end{table}

Step 6 is the one that decides, and the reason the rule ends with it is the
reason the whole book is organised the way it is. The accounting chain is
the reason the platform matters; the probability chain is the reason it
matters only sometimes. A rule that reversed the order would recommend
characterisation everywhere, and would be wrong.

\begin{keytakeaways}
\item The rule has six steps and ends with a test that can be failed: whether
the physical uncertainty dominates the encounter probability.
\item Applied to 2024 YR4's 2032 encounter, the rule declines the lunar
platform, because the astrometric term dominates by three orders of
magnitude.
\item Applied to Apophis's 2029 encounter, the rule selects it, because the
physical terms are each comparable with the astrometric standard deviation.
\item The rule returns a platform rather than a measurement, and it does not
carry a cost term the book cannot defend.
\end{keytakeaways}
